\section{PAT16 recovery: observation, uncertainty and known-value audit}\label{app:pat16_recovery}
\subsection{Fixed anatomy, initialization and physical observation region}
The experiment reuses the domain, lesion-aware tissues and transport law in Appendix~\ref{app:pat16}. Geometry and directions are fixed; the four migration/growth parameters are fitted. Initialization fixes $A=0.6$ and $\tau=0.5$ in the lesion smoother. The lesion-positive active-tissue centroid defines RAS center $(31.0259,35.7315,-12.6722)$ mm before signal generation. The fixed 72 mm cube has 12 centers per axis at 6 mm spacing, with full physical support inside the imaging grid: 1,728 observation voxels versus all 44,667 active PDE voxels. The prespecified 24 mm near-center subset contains 280 observation voxels.

For cube center $\bm x_v$, define offsets $\bm\delta_{abc}$ by a Cartesian product of $\{-2,0,2\}$ mm. The implemented imaging operator is
\begin{equation}
 H_v^h[u]=\frac{1}{27}\sum_{a,b,c=1}^{3}
 \mathcal I_h\!\left[G_{2\,\mathrm{mm}}*u\right](\bm x_v+\bm\delta_{abc}),\qquad
 y_{jv}=H_v^h[u(t_j)]+\epsilon_{jv},\quad
 \epsilon_{jv}\overset{\mathrm{iid}}{\sim}\mathcal N(0,0.015^2).
\end{equation}
Here $\mathcal I_h$ is physical-coordinate trilinear interpolation. The field is zero outside active tissue; Gaussian convolution uses constant-zero padding and interpolation uses nearest extension at image-center boundaries. Quadrature approximates a physical cube average. Post-operator noise is not clipped. Normalized signal scale, blur, averaging, noise SD and elapsed times are assumed known, not calibrated from preoperative MRI. Uniform $K$ cancels after setting $u=c/K$; recovery does not identify $K$ or absolute cell counts.

\subsection{Likelihood and numerical precision}
Write $\theta=(\log D_{\rm WM},\log D_{\rm GM},\log\rho,\log r)$, using the stated physical units as numerical references. The bounded Gaussian MLE minimizes
\begin{equation}
 \frac12\sum_{j:t_j\in\{20,40,60\}}\sum_v
 \left\{\frac{H_v^h[u_{\exp\theta}(t_j)]-y_{jv}}{0.015}\right\}^2.
\end{equation}
The bounds are $(0.05,0.02,0.01,1)$ to $(1,0.6,0.09,3.5)$ in physical parameter coordinates. Starts are $(0.25,0.20,0.025,1.5)$, $(0.40,0.15,0.035,2.0)$ and $(0.80,0.05,0.065,3.2)$; the smallest residual sum selects the reported fit. SciPy nonlinear least squares uses a three-point finite-difference Jacobian, relative step $10^{-4}$, Jacobian column scaling, $f_{\rm tol}=x_{\rm tol}=10^{-8}$ and $g_{\rm tol}=10^{-6}$. No prior or WM:GM constraint is used.

If the standardized Jacobian is $J=USV^{\mathsf T}$, all four singular values must exist and $s_{\min}>10^{-8}s_{\max}$. Otherwise fitting is rejected as locally unidentifiable. The log covariance is $V\operatorname{diag}(s_k^{-2})V^{\mathsf T}$, without estimating or rescaling the known noise variance. A marginal interval is $\exp[\hat\theta_k\pm1.96\sqrt{\Sigma_{kk}}]$. An optimum within $10^{-4}$ log units of either bound is flagged; such Wald intervals are not treated as calibrated. The inverse solver uses the pilot's diffusion/logistic update and CFL subdivision, but retains float64 snapshots to avoid rounding error in likelihood derivatives. Tests check agreement with the pilot's float32 output and double-precision mass conservation.

\subsection{Resampling, validation and forecast intervals}
Seed 161026 spawns four independent NumPy streams for primary scan noise, bootstrap data, repeated-noise data and predictive observation noise. Changing replication counts therefore preserves the primary realization and the other streams. Bootstrap early scans are generated at the primary fitted parameters, with \PRBootstraps{} refits initialized at that fit. Percentile intervals use successful draws, including boundary draws if present; every failure and boundary flag is retained and no replacement draw is generated.

For the \PRTrials{}-trial known-value audit, fresh independent noise is added to the 0.25-day generator's early signals at the known parameters. Each refit starts from the primary estimate. Mean, relative bias and RMSE use successful fits; marginal local coverage counts every attempted trial and treats failed/boundary intervals as misses. Exact binomial intervals describe Monte Carlo uncertainty in each coverage fraction. None failed or reached a bound. Neither marginal percentile-bootstrap parameter-interval coverage nor joint four-parameter coverage is audited.

Each bootstrap vector predicts all observation times without consulting days 75/90. The 2.5th/97.5th percentiles give nominal pointwise 95\% signal intervals. The original prediction construction added one independent Gaussian noise field to each of 100 bootstrap trajectories, introducing finite-tail sampling error. For the revised intervals, let $\mu_{bv}$ denote the forecast mean at a fixed time and voxel for bootstrap vector $b$. We integrate the known observation noise:
\begin{equation}
 F_v(q)=\frac{1}{B}\sum_{b=1}^{B}\Phi\!\left(\frac{q-\mu_{bv}}{0.015}\right),\qquad
 [L_v,U_v]=[F_v^{-1}(0.025),F_v^{-1}(0.975)],\quad B=100.
\end{equation}
Quantiles use 32 bisection iterations within the smallest/largest component quantiles. These are quantiles of an empirical bootstrap/Gaussian mixture, not exact frequentist coverage bounds. For comparison, 1,000 and \PPDraws{} draws resample a whole bootstrap trajectory index and add independent noise per voxel/time. Spatial chunks of 128 points bound memory and use fixed seeds; the trajectory index is shared across chunks. The 20,000 draws improve noise-tail precision but still reuse the same 100 parameter vectors.

A separate seed 161027 specifies \PPTrials{} new complete datasets. Trial index and stream identifiers independently seed training noise, bootstrap noise, future test noise, legacy predictive noise and 20,000-draw predictive resampling. Each trial generates days 20/40/60 from the fixed 0.25-day known signal, refits from $(0.4,0.15,0.035,2)$, generates and refits 100 early bootstrap datasets at that trial's fitted model, and forecasts days 75/90. The same bootstrap forecasts compare legacy, 20,000-draw and integrated intervals, paired within trial. Future observations are generated independently at the known finer-time signal. No failed fit is replaced. Unusable primary fits (failure/boundary or fewer than two successful bootstrap refits) count as coverage misses; widths/scores would be reported only for usable trials. Boundary bootstrap draws are retained and flagged. All \PPTrials{} primary and \PPNestedFits{} bootstrap fits succeeded without reaching a bound.

For trial $i$, region $R$ and day $t$, let $c_i=|R|^{-1}\sum_{v\in R}\mathbf1\{L_{iv}\le y_{iv}\le U_{iv}\}$. We report $\bar c$ with Monte Carlo interval $\bar c\pm t_{0.975,M-1}s_c/\sqrt M$, $M=\PPTrials{}$. This measures uncertainty in region-averaged marginal coverage across independent complete experiments; within-image voxels share fitted parameters and are not pooled as binomial replicates. The fixed central index $(6,6,6)$ (zero-based; inside the prespecified sphere) additionally has exact binomial coverage intervals over the $M$ independent future observations. Mean width and the proper 95\% interval score \citep{gneiting2007}
\begin{equation}
 S(L,U;y)=U-L+\frac{2}{0.05}(L-y)\mathbf1\{y<L\}
                   +\frac{2}{0.05}(y-U)\mathbf1\{y>U\}
\end{equation}
use the same trials and regions. Records also retain conditional expected containment from the known Gaussian generator and paired changes from the legacy procedure. Neither is used to tune the interval. These checks validate only this fixed design and noise model; pointwise coverage does not imply simultaneous image coverage.

Two additional noise-free fits isolate recoverability and temporal sensitivity. The inverse model's own signal must recover known parameters within maximum relative error $10^{-4}$. Fitting noise-free signals from the finer-time generator then quantifies deterministic temporal mismatch. Both retain the same spatial anatomy, initialization, PDE family and imaging operator; they do not independently validate spatial convergence or biology. A time-zero-only fit must fail the rank criterion. Separate tests remove directional eligibility or one tissue class to demonstrate the corresponding unidentifiable parameters.

\subsection{Public records and manuscript generation}
\path{pat16_synthetic_parameter_recovery.ipynb} replays checked records; refitting requires real derivatives through \path{PAT16_SOURCE}. \path{scripts/run_pat16_recovery.py} runs the audit. \path{outputs/pat16_recovery/} contains the summary, parameter CSV, observation/resampling arrays and figures, retaining settings, failures, forecasts and numerical checks. Provenance records input, source and output hashes and software versions; raw MRI is not copied. The pilot loader requires its full derivative set, although fMRI is absent from the likelihood.

\path{paper/make_pat16_recovery_figures.py} verifies hashes, the selected multistart fit and both primary/predictive noise streams. It recomputes intervals, bias/RMSE/coverage, binomial limits, regions, forecast errors/widths/containment and numerical discrepancies. It generates three figures, three recovery tables, numerical macros and \path{paper/data/pat16_recovery_claims_audit.json} without MRI or refitting. The additional \path{src/brainweb_kpp/pat16_prediction.py} and \path{scripts/run_pat16_prediction.py} implement the nested predictive audit. \path{outputs/pat16_prediction/} saves the primary comparisons and a separate checkpoint for every complete trial, retaining training/test data, bootstrap vectors/forecasts, all three interval sets, metrics and failure flags. Its manifest hashes source, baseline and output files; checkpoints permit deterministic resumption with the same configuration. \path{paper/make_pat16_prediction_figures.py} verifies hashes, independent noise streams, Gaussian-mixture CDF endpoints, containment, widths, scores and coverage uncertainty before generating the additional validation table/figure and claims audit. Default paper builds replay both audits without MRI. \texttt{make pat16-prediction} reruns the nested experiment with eight independent worker processes. \texttt{make pat16-recovery} reruns the MRI-dependent audit. Sampling intervals exclude geometry, registration, segmentation, directions, initialization, signal calibration and shared spatial/model discrepancy. No patient rate estimate, functional forecast or clinical validation is claimed.
